\chapter*{Danh mục thuật ngữ và từ viết tắt}

Danh mục này giải thích các thuật ngữ chuyên môn và từ viết tắt có vai trò quan trọng trong báo cáo ConferenceSpace. Các mục được chọn theo mức độ xuất hiện và vai trò diễn giải trong Chương~1--5; tên tiếng Anh chuẩn được giữ khi báo cáo dùng chúng như nhãn kỹ thuật hoặc chỉ số, còn phần tiếng Việt bám các cụm đã thống nhất trong toàn báo cáo.

\section*{Thuật ngữ nghiệp vụ hội nghị}

\begingroup
\small
\setlength{\tabcolsep}{2pt}
\renewcommand{\arraystretch}{1.2}
\begin{longtable}{|C{0.22\textwidth}|P{0.26\textwidth}|P{0.46\textwidth}|}
  \caption{Danh mục thuật ngữ và từ viết tắt, phần 1}\\
  \hline
   \TableHeader{Thuật ngữ}    &    \TableHeader{Tên đầy đủ / tương đương}    &    \TableHeader{Nghĩa sử dụng trong báo cáo} \\ \hline
  Chair     &     Chủ tọa (Chair)     &     Vai trò chịu trách nhiệm cấu hình hội nghị, theo dõi tiến độ, phân công phản biện và ra quyết định cuối cùng đối với bài nộp. Trong ConferenceSpace thường viết tắt là \textbf{Chủ tọa}; khi cần phân biệt với vai trò đồng điều phối dùng \textbf{Chủ tọa/Đồng chủ tọa} \\ \hline
  Area Chair      &      Chủ tọa khu vực (Area Chair)      &      Vai trò điều phối theo khu vực chủ đề trong hội nghị quy mô lớn (ví dụ OpenReview/NeurIPS); xuất hiện khi báo cáo đối chiếu quy trình hội nghị thực tế, không phải vai trò nghiệp vụ chính của ConferenceSpace \\ \hline
  Reviewer      &      phản biện viên      &      Người được phân công đọc bài, viết \textbf{nhận xét / bản phản biện}, chấm điểm hoặc đưa ra khuyến nghị trong quy trình xét duyệt. \textbf{phản biện viên} là vai trò; \textbf{nhận xét phản biện} / \textbf{bản phản biện} là sản phẩm do phản biện viên viết \\ \hline
  Submission     &     bài nộp     &     Bản thảo và metadata do tác giả gửi vào một hội nghị hoặc chuyên đề (track) cụ thể. Giữ tiếng Anh trong tên luồng sản phẩm (ví dụ Submission Autofill, Submission Gating) và trong mã kỹ thuật \\ \hline
  CFP      &      Call for Papers      &      Lời mời nộp bài của hội nghị, thường mô tả phạm vi chủ đề, chuyên đề, thời hạn và yêu cầu nộp bài \\ \hline
  Rebuttal     &     giai đoạn phản hồi của tác giả (rebuttal)     &     Giai đoạn tác giả trả lời nhận xét của phản biện viên trước khi Chủ tọa ra quyết định cuối cùng. Sau lần giới thiệu đầy đủ, báo cáo có thể rút gọn thành \textit{phản hồi của tác giả (rebuttal)} \\ \hline
  Camera-ready      &      bản camera-ready / bản hoàn thiện sau chấp nhận      &      Phiên bản cuối của bài báo sau khi bài được chấp nhận và tác giả đã chỉnh sửa theo yêu cầu hội nghị \\ \hline
  Bidding     &     đăng ký/chọn bài phản biện (bidding)     &     Cơ chế phản biện viên bày tỏ mức độ quan tâm hoặc phù hợp với các bài nộp trước khi phân công. Trong phạm vi hiện tại, đây là hướng hoàn thiện nghiệp vụ, chưa phải luồng đã đánh giá đầy đủ \\ \hline
  Peer review      &      bình duyệt học thuật (peer review)      &      Quy trình chuyên gia đánh giá chất lượng, tính mới, độ đúng và mức phù hợp của bài báo khoa học; còn gọi là phản biện học thuật khi nhấn mạnh hoạt động phản biện \\ \hline
\end{longtable}
\endgroup

\section*{Thuật ngữ hệ thống và AI}

\begingroup
\small
\setlength{\tabcolsep}{2pt}
\renewcommand{\arraystretch}{1.2}
\begin{longtable}{|C{0.22\textwidth}|P{0.26\textwidth}|P{0.46\textwidth}|}
  \caption{Danh mục thuật ngữ và từ viết tắt, phần 2}\\
  \hline
   \TableHeader{Thuật ngữ}    &    \TableHeader{Tên đầy đủ / tương đương}    &    \TableHeader{Nghĩa sử dụng trong báo cáo} \\ \hline
  AI     &     Artificial Intelligence     &     Nhóm chức năng hỗ trợ bằng mô hình trí tuệ nhân tạo trong hệ thống. Không bao gồm các cơ chế thuật toán cố định, có thể kiểm chứng như đối sánh phản biện (reviewer matching) và phát hiện xung đột lợi ích (COI) \\ \hline
  LLM     &     Large Language Model     &     Mô hình ngôn ngữ lớn dùng trong các luồng hỗ trợ nhập liệu, đọc bài, kiểm tra bản nháp nhận xét phản biện, tổng hợp bằng chứng hoặc chatbot \\ \hline
  Luồng AI      &      luồng xử lý AI / AI workflow      &      Luồng xử lý có dùng AI, có đầu vào, đầu ra, ranh giới kiểm soát và tiêu chí đánh giá riêng. Báo cáo ưu tiên từ \textbf{luồng} thay cho lặp \textit{workflow} khi không cần tên riêng tiếng Anh \\ \hline
  Sáu luồng AI chính     &     tên sản phẩm tiếng Anh giữ nguyên     &     Submission Autofill (tự động điền bài nộp); Submission Gating (kiểm tra trước khi nộp bài); Reviewer Initial Analysis (phân tích và tóm tắt bài báo); Review Quality Auditor (kiểm tra chất lượng nhận xét phản biện); Chair Decision Copilot (báo cáo tổng hợp hỗ trợ Chủ tọa); Chatbot Agent (trợ lý hội thoại có công cụ và phân quyền) \\ \hline
  Thuật toán cố định, có thể kiểm chứng      &      deterministic / thuật toán xác định      &      Lớp xử lý đối sánh phản biện và COI: kết quả ổn định, có thể tái tạo và kiểm tra lại; không dùng AI tạo sinh. Lần đầu trong đoạn dùng cụm đầy đủ; các lần sau có thể xen kẽ \textit{thuật toán có thể kiểm chứng} hoặc \textit{thuật toán xác định} \\ \hline
  Agent      &      tác nhân (agent)      &      Thành phần AI có thể chọn chiến lược truy xuất, chia nhỏ tác vụ, gọi công cụ và lặp bước kiểm tra trong phạm vi quyền. Không dùng \textit{tác tử} \\ \hline
  human-in-the-loop     &     sự can thiệp của con người (human-in-the-loop)     &     Bước người hoặc vai trò có thẩm quyền xem lại, chỉnh sửa, ghi đè hoặc xác nhận trước khi dùng đầu ra AI gần quyết định học thuật. Không dùng \textit{vòng kiểm soát của con người} \\ \hline
  ReAct     &     Reasoning and Acting     &     Mẫu kết hợp suy luận và gọi công cụ cho agent nhiều bước. Báo cáo chỉ giữ một lớp gloss tiếng Anh, không nhét thêm bản dịch tiếng Việt kép \\ \hline
  Bộ chạy luồng xử lý     &     workflow runner     &     Thành phần chạy luồng AI trên tập dữ liệu benchmark, lưu đầu ra, điểm kiểm tra, thời gian xử lý, token và trạng thái hoàn tất \\ \hline
  Dispatcher-worker     &     bộ điều phối và worker xử lý     &     Mô hình vận hành trong đó dispatcher phân phối tác vụ AI cho các worker chạy độc lập và lưu kết quả theo từng tác vụ \\ \hline
  Artifact      &      hiện vật / gói kết quả kỹ thuật      &      Kết quả hoặc đối tượng kỹ thuật được tạo và lưu để đối chiếu, ví dụ đầu ra của luồng AI, metadata kiểm toán hoặc dữ liệu đã gắn fingerprint \\ \hline
  Structured output     &     đầu ra có cấu trúc (structured output)     &     Đầu ra của mô hình phải tuân theo schema hoặc ràng buộc field để backend có thể kiểm tra, lưu và hiển thị \\ \hline
  Schema      &      lược đồ dữ liệu      &      Đặc tả cấu trúc dữ liệu, kiểu field và ràng buộc kiểm tra cho request/response, đầu ra AI hoặc artifact \\ \hline
  data contract      &      ràng buộc dữ liệu (data contract) / chuẩn đầu ra      &      Quy ước bắt buộc về cấu trúc và điều kiện đầu ra giữa AI service và backend. Không dùng calque \textit{hợp đồng dữ liệu} / \textit{hợp đồng đầu ra}; khi nhấn mạnh hành động kiểm tra dùng \textit{đối chiếu đầu ra với chuẩn đã định sẵn} \\ \hline
  Composite pattern      &      mẫu Composite      &      Mẫu thiết kế gom nhiều detector hoặc thành phần xử lý cùng giao diện để bật/tắt hoặc kết hợp theo cấu hình, dùng trong phát hiện COI nhiều lớp \\ \hline
  Hallucination     &     ảo tưởng (hallucination) / thông tin thiếu căn cứ     &     Trường hợp AI tạo nội dung không được dữ liệu đầu vào hoặc bằng chứng hệ thống hỗ trợ. Báo cáo cũng diễn đạt là \textit{thông tin thiếu căn cứ} hoặc \textit{nội dung do AI tạo ra không có căn cứ} \\ \hline
  Retrieval-Augmented Generation     &     truy hồi tăng cường tạo sinh     &     Cách kết hợp mô hình ngôn ngữ với bước truy xuất tài liệu hoặc artifact liên quan trước khi sinh câu trả lời, nhằm tăng mức bám nguồn \\ \hline
  RAG      &      Retrieval-Augmented Generation      &      Viết tắt của Retrieval-Augmented Generation; trong báo cáo chủ yếu là hướng phát triển để tăng groundedness của luồng AI \\ \hline
  Agentic RAG     &     truy hồi tăng cường tạo sinh dạng tác nhân (Agentic RAG)     &     Hướng mở rộng trong đó agent chọn chiến lược truy xuất, chia nhỏ tác vụ, gọi công cụ và kiểm tra lại nguồn thay vì chạy một pipeline cố định \\ \hline
  Vector search      &      tìm kiếm vector (vector search)      &      Cách tìm tài liệu, chủ đề hoặc hồ sơ gần nhau trong không gian embedding thay vì chỉ khớp từ khóa \\ \hline
  Hybrid retrieval     &     truy hồi lai (hybrid retrieval)     &     Cách kết hợp tìm kiếm từ khóa, metadata và vector search để tăng khả năng lấy đúng nguồn liên quan \\ \hline
  Observability      &      khả năng quan sát (observability)      &      Khả năng theo dõi luồng xử lý đã dùng nguồn nào, gọi công cụ nào, sinh đầu ra nào, lỗi ở đâu và người dùng đã chấp nhận, chỉnh sửa hoặc ghi đè kết quả ra sao \\ \hline
  AgentOps     &     vận hành tác nhân (AgentOps)     &     Nhóm thực hành giám sát, ghi log, theo dõi vết, analytics và kiểm soát an toàn cho hệ thống agent, tương tự DevOps/MLOps cho agent \\ \hline
  Audit trail      &      dấu vết kiểm toán (audit trail)      &      Chuỗi sự kiện và artifact giúp truy lại thành phần nào đã thực hiện thao tác, dựa trên dữ liệu nào và tạo kết quả gì \\ \hline
  Rollback      &      khôi phục (rollback)      &      Cơ chế đưa hệ thống hoặc artifact về trạng thái trước đó khi thao tác tự động hoặc bán tự động không phù hợp \\ \hline
  Fingerprint     &     dấu vết trạng thái / mã định danh (fingerprint)     &     Giá trị đại diện cho trạng thái dữ liệu tại thời điểm sinh artifact, dùng để phát hiện artifact đã lỗi thời \\ \hline
  Stale      &      không còn hiện hành (stale)      &      Trạng thái của artifact khi dữ liệu gốc đã thay đổi so với thời điểm artifact được tạo \\ \hline
  Model router     &     bộ định tuyến mô hình (model router)     &     Lớp chọn hoặc gọi nhà cung cấp mô hình theo cấu hình runtime của AI service \\ \hline
  Open scholarly infrastructure      &      hạ tầng học thuật mở (open scholarly infrastructure)      &      Định hướng vận hành hạ tầng phục vụ cộng đồng học thuật với quản trị, tính bền vững, mã nguồn mở hoặc chuẩn mở phù hợp; xuất hiện khi bàn mô hình vận hành dài hạn \\ \hline
  PoC     &     mô hình kiểm chứng (PoC / Proof of Concept)     &     Bản chứng minh khả thi về mặt kỹ thuật và học thuật, chưa nhằm chứng minh mô hình kinh doanh hoặc khả năng vận hành thương mại đầy đủ \\ \hline
\end{longtable}
\endgroup

\section*{Thuật ngữ phân quyền và Chatbot Agent}

\begingroup
\small
\setlength{\tabcolsep}{2pt}
\renewcommand{\arraystretch}{1.2}
\begin{longtable}{|C{0.22\textwidth}|P{0.26\textwidth}|P{0.46\textwidth}|}
  \caption{Danh mục thuật ngữ và từ viết tắt, phần 3}\\
  \hline
   \TableHeader{Thuật ngữ}    &    \TableHeader{Tên đầy đủ / tương đương}    &    \TableHeader{Nghĩa sử dụng trong báo cáo} \\ \hline
  API     &     Application Programming Interface     &     Giao diện cho phép các thành phần hệ thống trao đổi request/response \\ \hline
  RBAC     &     Role-Based Access Control     &     Cơ chế kiểm soát quyền truy cập dựa trên vai trò của người dùng trong từng hội nghị và trên từng tài nguyên \\ \hline
  Resource registry      &      danh mục tài nguyên (resource registry)      &      Danh sách tài nguyên mà agent được phép mô tả và truy vấn, kèm kiểm tra field và chính sách, ví dụ bài nộp, phân công hoặc thông báo \\ \hline
  Query engine      &      backend query engine      &      Công cụ truy vấn chỉ đọc để Chatbot Agent lấy dữ liệu hệ thống thông qua backend và lớp phân quyền, không truy cập cơ sở dữ liệu trực tiếp \\ \hline
  Service-to-service token      &      token nội bộ giữa các service      &      Token xác thực dùng cho giao tiếp nội bộ giữa backend và AI service trong phạm vi đã kiểm soát \\ \hline
  Ranh giới phân quyền     &     permission boundary / ranh giới quyền     &     Giới hạn dữ liệu hoặc thao tác mà một vai trò người dùng được phép truy cập; đặc biệt quan trọng với Chatbot Agent để tránh rò rỉ bản thảo hoặc nhận xét phản biện \\ \hline
  Khả năng tuân thủ quyền truy cập     &     permission compliance     &     Mức độ chatbot hoặc agent giữ đúng ranh giới quyền, không trả dữ liệu vượt quyền của vai trò đang sử dụng. Thay cho nhãn \textit{permission safety} không dùng trong báo cáo \\ \hline
  SLA      &      Service Level Agreement      &      Cam kết mức dịch vụ, ví dụ độ sẵn sàng, thời gian phản hồi hoặc hỗ trợ vận hành trong mô hình dịch vụ được vận hành (hosted service) \\ \hline
\end{longtable}
\endgroup

\section*{Thuật ngữ đối sánh phản biện và COI}

\begingroup
\small
\setlength{\tabcolsep}{2pt}
\renewcommand{\arraystretch}{1.2}
\begin{longtable}{|C{0.22\textwidth}|P{0.26\textwidth}|P{0.46\textwidth}|}
  \caption{Danh mục thuật ngữ và từ viết tắt, phần 4}\\
  \hline
   \TableHeader{Thuật ngữ}    &    \TableHeader{Tên đầy đủ / tương đương}    &    \TableHeader{Nghĩa sử dụng trong báo cáo} \\ \hline
  COI     &     Conflict of Interest / xung đột lợi ích     &     Xung đột lợi ích giữa phản biện viên và bài nộp, cần được phát hiện trước khi phân công phản biện \\ \hline
  Reviewer matching      &      đối sánh phản biện (reviewer matching)      &      Cơ chế gợi ý phản biện viên phù hợp với bài nộp dựa trên miền chuyên môn, khối lượng công việc và ràng buộc COI. Lớp này thuộc thuật toán cố định, có thể kiểm chứng, không phải luồng AI tạo sinh \\ \hline
  Domain Jaccard     &     Jaccard theo miền chuyên môn     &     Công thức đo độ tương đồng giữa tập từ khóa/miền của bài nộp và phản biện viên bằng tỷ lệ giao trên hợp \\ \hline
  Greedy assignment      &      phân công tham lam (greedy)      &      Thuật toán xét các cặp bài nộp--phản biện viên có điểm cao trước, đồng thời tôn trọng giới hạn tải và COI \\ \hline
  Cân bằng tải     &     load balancing     &     Nguyên tắc phân bổ phân công để tránh dồn quá nhiều bài cho một phản biện viên khi có nhiều ứng viên tương đương \\ \hline
  Độ lệch chuẩn tải      &      load standard deviation      &      Chỉ số đo mức phân tán số bài được gán cho từng phản biện viên; giá trị cao cho thấy tải phân công lệch hơn. Báo cáo không dùng nhãn \textit{Load StdDev} \\ \hline
  Hệ số Gini tải     &     load Gini coefficient     &     Chỉ số đo mức bất bình đẳng trong phân bổ bài cho phản biện viên; giá trị càng gần 0 thì tải càng đều. Báo cáo dùng \textit{hệ số Gini tải}, không dùng nhãn \textit{Load Gini} \\ \hline
  Tie-break      &      quy tắc phá hòa      &      Quy tắc chọn thứ tự ưu tiên khi nhiều ứng viên có cùng điểm, ví dụ ưu tiên phản biện viên có tải thấp hơn \\ \hline
  Fallback     &     bước dự phòng (fallback)     &     Bước xử lý khi một bài chưa đủ phản biện viên sau lượt gán chính; có thể nới điểm hoặc giới hạn tải nhưng không nới COI \\ \hline
  Tỷ lệ chuyển sang phân công ngẫu nhiên      &      fallback / tỷ lệ thoái lui (khi diễn đạt ngắn)      &      Tỷ lệ phân công phải dùng cơ chế dự phòng vì không còn phản biện viên hợp lệ có tín hiệu phù hợp đủ rõ. Cụm chính trong Chương~4 là \textit{tỷ lệ chuyển sang phân công ngẫu nhiên}; không dùng nhãn \textit{fallback rate} \\ \hline
  Vi phạm COI     &     vi phạm xung đột lợi ích (COI)     &     Trường hợp hệ thống phân công phản biện viên vào bài nộp có COI; trong benchmark đây là chỉ số phải bằng 0 \\ \hline
  Co-author graph      &      đồ thị đồng tác giả      &      Đồ thị biểu diễn quan hệ đồng tác giả giữa các nhà nghiên cứu, dùng để phát hiện COI gián tiếp trên Neo4j \\ \hline
  Trần từ vựng     &     lexical ceiling     &     Giới hạn của phương pháp so khớp từ khóa khi hai hồ sơ cùng lĩnh vực lớn nhưng dùng ít từ khóa trùng nhau \\ \hline
  Embedding     &     biểu diễn nhúng / vector ngữ nghĩa     &     Cách biểu diễn văn bản hoặc chủ đề thành vector để đo tương đồng ngữ nghĩa sâu hơn so với trùng khớp từ khóa \\ \hline
  ANN      &      Approximate Nearest Neighbor      &      Kỹ thuật tìm láng giềng gần xấp xỉ trong không gian vector, thường dùng khi cần xếp hạng tương đồng trên tập dữ liệu lớn \\ \hline
\end{longtable}
\endgroup

\section*{Thuật ngữ đánh giá thực nghiệm}

\begingroup
\small
\setlength{\tabcolsep}{2pt}
\renewcommand{\arraystretch}{1.2}
\begin{longtable}{|C{0.22\textwidth}|P{0.26\textwidth}|P{0.46\textwidth}|}
  \caption{Danh mục thuật ngữ và từ viết tắt, phần 5}\\
  \hline
   \TableHeader{Thuật ngữ}    &    \TableHeader{Tên đầy đủ / tương đương}    &    \TableHeader{Nghĩa sử dụng trong báo cáo} \\ \hline
  Benchmark     &     thực nghiệm đánh giá có kịch bản     &     Quy trình đo hiệu năng, chất lượng hoặc độ tin cậy của một lớp hệ thống theo dữ liệu và chỉ số đã xác định \\ \hline
  Go micro-benchmark      &      (giữ tên tiếng Anh)      &      Thực nghiệm đo trực tiếp thời gian xử lý và mức sử dụng bộ nhớ của thuật toán trong chương trình, \textbf{không bao gồm} chi phí của HTTP, cơ sở dữ liệu và mạng \\ \hline
  k6     &     công cụ kiểm thử tải HTTP     &     Công cụ dùng để sinh tải request trong benchmark backend, đo độ trễ, thông lượng và tỷ lệ lỗi \\ \hline
  Người dùng ảo      &      virtual user      &      Đơn vị mô phỏng người dùng đồng thời trong kiểm thử tải bằng k6 \\ \hline
  Đường cơ sở     &     baseline     &     Phương pháp hoặc kết quả nền dùng để so sánh, ví dụ xếp hạng ngẫu nhiên trong đánh giá đối sánh phản biện \\ \hline
  Snapshot      &      bản sao lưu dữ liệu (snapshot)      &      Bản dữ liệu Semantic Scholar hoặc tập dữ liệu cố định dùng khi đánh giá, không gọi API lúc chạy thử nghiệm. Không dùng \textit{ảnh chụp dữ liệu} \\ \hline
  Dữ liệu tham chiếu     &     nhãn tham chiếu / ground truth     &     Dữ liệu hoặc nhãn được xem là chuẩn đối chiếu trong một thực nghiệm \\ \hline
  Proxy      &      chỉ số gián tiếp (proxy)      &      Chỉ số gián tiếp dùng khi không thể đo trực tiếp chất lượng cần quan tâm, ví dụ dùng TCA để đánh giá mức bám chứng cứ của đầu ra AI. Không dùng \textit{chỉ báo thay thế} hay \textit{bằng chứng thay thế} \\ \hline
  TCA     &     Textual Claim-based Assessment     &     Khung đánh giá đầu ra AI theo mệnh đề--chứng cứ, gồm độ trung thực (Truthfulness), độ phủ (Coverage) và khả năng bổ sung (Additionality). Không expand nhầm thành chuỗi ghép ba từ tiếng Anh \\ \hline
  NLI      &      Natural Language Inference / suy luận ngôn ngữ tự nhiên      &      Kỹ thuật kiểm tra quan hệ suy luận giữa mệnh đề và chứng cứ, dùng trong benchmark đánh giá một số đầu ra AI \\ \hline
  Truthfulness     &     độ trung thực / mức bám nguồn     &     Tỷ lệ hoặc mức độ các mệnh đề được nguồn dữ liệu hỗ trợ \\ \hline
  Groundedness      &      mức độ bám chứng cứ (groundedness)      &      Mức độ một nhận định có thể truy vết về chứng cứ cụ thể trong dữ liệu đầu vào \\ \hline
  Coverage (TCA)     &     độ phủ (Coverage)     &     Trong benchmark AI: mức độ đầu ra bao phủ các điểm có trong nguồn tham chiếu hoặc trong nội dung do con người tạo ra \\ \hline
  Độ phủ phân công      &      assignment coverage      &      Trong đối sánh phản biện: tỷ lệ bài được phân công đủ phản biện viên theo yêu cầu thử nghiệm \\ \hline
  Additionality      &      khả năng bổ sung (Additionality)      &      Phần nội dung đúng và có chứng cứ nhưng không trùng trực tiếp với tham chiếu của con người \\ \hline
  Tỷ lệ vừa bám nguồn vừa hợp lệ     &     grounded-and-valid rate     &     Chỉ số kết hợp giữa việc cảnh báo hoặc phát hiện có căn cứ và việc phù hợp với tiêu chí kiểm tra. Thay cho nhãn \textit{Grounded-valid rate} \\ \hline
  Validity     &     tính hợp lệ / validity rate     &     Tỷ lệ phát hiện hoặc cảnh báo được đánh giá là phù hợp với tiêu chí kiểm tra, chưa nhất thiết đảm bảo nội dung đó bám nguồn đầy đủ \\ \hline
  Mức độ nghiêm trọng     &     severity     &     Nhãn thể hiện độ nặng của cảnh báo nội dung hoặc phát hiện, ví dụ warn/block hoặc minor/moderate/major \\ \hline
  Tỷ lệ hoàn tất     &     completion rate     &     Tỷ lệ tác vụ, luồng xử lý hoặc hội thoại chạy xong và tạo được đầu ra hợp lệ theo chuẩn đã định sẵn \\ \hline
  Tỷ lệ chuyên đề không hợp lệ      &      invalid track rate      &      Tỷ lệ gợi ý chuyên đề (track) nằm ngoài danh sách chuyên đề hợp lệ của hội nghị đang xét \\ \hline
  Tỷ lệ rủi ro cao     &     high-risk rate     &     Tỷ lệ đầu ra hoặc tác vụ có dấu hiệu cần kiểm tra thủ công kỹ hơn trước khi dùng gần điểm quyết định nhạy cảm \\ \hline
  Exact Match      &      tỉ lệ khớp chính xác (Exact Match)      &      Chỉ số yêu cầu đầu ra trùng hoàn toàn với giá trị tham chiếu sau khi chuẩn hóa theo quy tắc đánh giá \\ \hline
  Hit@K      &      Hit at K      &      Chỉ số xếp hạng cho biết kết quả đúng có nằm trong $K$ gợi ý đầu tiên hay không. Báo cáo cũng viết Hit@k \\ \hline
  MRR     &     Mean Reciprocal Rank     &     Trung bình nghịch đảo thứ hạng của kết quả đúng đầu tiên trong danh sách gợi ý \\ \hline
  nDCG@K      &      Normalized Discounted Cumulative Gain at K      &      Chỉ số đánh giá xếp hạng có xét vị trí và độ liên quan của các kết quả trong $K$ vị trí đầu \\ \hline
  F1     &     F1-score     &     Trung bình điều hòa giữa precision và recall, dùng trong các bài toán so khớp hoặc trích xuất metadata \\ \hline
  ROUGE      &      Recall-Oriented Understudy for Gisting Evaluation      &      Nhóm chỉ số so sánh độ trùng lặp văn bản, dùng trong đánh giá tóm tắt hoặc trích xuất abstract \\ \hline
  p95 latency     &     độ trễ phân vị 95     &     Mức độ trễ mà 95\% request có thời gian xử lý nhỏ hơn hoặc bằng giá trị này \\ \hline
  Throughput      &      thông lượng      &      Số request hoặc tác vụ hệ thống xử lý được trong một đơn vị thời gian \\ \hline
  Token     &     đơn vị văn bản của mô hình     &     Đơn vị dùng để đo lượng đầu vào/đầu ra và ước tính chi phí khi gọi LLM. Khác với token xác thực (JWT, service token) \\ \hline
  TTFT      &      thời gian đến token đầu tiên (TTFT)      &      Thời gian từ lúc gửi yêu cầu đến khi hệ thống bắt đầu trả token đầu tiên trong luồng phản hồi chatbot \\ \hline
  Thời lượng luồng phản hồi     &     stream duration     &     Thời gian hệ thống duy trì luồng trả lời sau khi bắt đầu stream dữ liệu về giao diện. Thay cho nhãn \textit{Stream duration} đứng một mình \\ \hline
  Tỷ lệ gọi công cụ thành công      &      tool-call success rate      &      Tỷ lệ lời gọi công cụ (tool call) của agent trả về kết quả hợp lệ trong benchmark hội thoại \\ \hline
  Kết quả hội thoại     &     conversation / workflow outcome     &     Nhãn đánh giá cuối cùng của một hội thoại chatbot, gồm đạt, đạt một phần hoặc không đạt \\ \hline
  UAT     &     User Acceptance Testing / khảo sát sau sử dụng     &     Khảo sát hoặc kiểm thử mức chấp nhận của người dùng sau khi trải nghiệm hệ thống, trên các vai trò Tác giả, phản biện viên và Chủ tọa \\ \hline
\end{longtable}
\endgroup

\section*{Từ viết tắt triển khai}

\begingroup
\small
\setlength{\tabcolsep}{2pt}
\renewcommand{\arraystretch}{1.2}
\begin{longtable}{|C{0.22\textwidth}|P{0.26\textwidth}|P{0.46\textwidth}|}
  \caption{Danh mục thuật ngữ và từ viết tắt, phần 6}\\
  \hline
   \TableHeader{Thuật ngữ}    &    \TableHeader{Tên đầy đủ / tương đương}    &    \TableHeader{Nghĩa sử dụng trong báo cáo} \\ \hline
  CI/CD     &     Continuous Integration / Continuous Deployment     &     Quy trình tự động build, kiểm tra, đóng gói và triển khai hệ thống \\ \hline
  GHCR      &      GitHub Container Registry      &      Registry lưu container image được build từ GitHub Actions \\ \hline
  VPS     &     Virtual Private Server     &     Máy chủ ảo dùng để triển khai môi trường production trong phạm vi đề tài \\ \hline
  HTTP      &      Hypertext Transfer Protocol      &      Giao thức request/response dùng cho API và benchmark tải backend \\ \hline
  HTTPS     &     HTTP Secure     &     HTTP qua kênh mã hóa TLS, được Caddy xử lý trong triển khai production \\ \hline
  CRUD      &      Create, Read, Update, Delete      &      Nhóm thao tác nghiệp vụ cơ bản trên dữ liệu hệ thống; cũng là một kịch bản tải HTTP trong Chương~4 \\ \hline
\end{longtable}
\endgroup
